\subsection{Reproduction on the archimedean frame slab}
\label{arrep:section}

Write $\Psi$ for the CM type, $p=p_0$, $q=r-p$, and $G=G_\tau$
for the Hermitian matrix at $\tau$. The pairing is linear in its
first variable, so the Gram matrix of a column frame $x$ is
$x^tG\bar x$. Put $J=\operatorname{diag}(1_p,-1_q)$. The domain
$\mathcal D$ consists of invertible frames at all $\tau$ whose first $p$
columns span a positive space. This is equivalent to the prescribed
signs of all the leading Gram determinants: after the first $p$
columns, the additional orthogonal directions belong to a negative
definite space. Choose an orthoframe $x_0$ with Gram matrix $J$ at
every place. Let $B$ be the product of the invertible upper triangular
groups, and let $dy$ be ordinary Lebesgue measure in the original
complex matrix coordinates. Set
\[
 S_\delta=\{y:\ |(y_\tau^tG_\tau\bar y_\tau-J)_{ij}|<\delta
                         \text{ for every }\tau,i,j\}.
\]

For each place choose nonnegative integers $c_1,\ldots,c_q$,
$c'_1,\ldots,c'_p$ and write $a_i=\sum_{j\ge i}c_j$,
$a'_i=\sum_{j\ge i}c'_j$. In physical column order the character
$\chi_\lambda(b)=\prod_{\tau,i}b_{\tau,ii}^{\lambda_{\tau,i}}$ has weight
\[
 \lambda_\tau=(-n/2-a'_p,\ldots,-n/2-a'_1,\;
                          n/2+a_1,\ldots,n/2+a_q).
\]
The polynomial letters at each place have $n$ columns and two row
blocks. The $q$ rows are the negative rows in their physical order,
and the $p$ rows are the positive rows in reverse physical order.
A row-highest harmonic polynomial has row degrees $a_i,a'_i$, is
invariant under lower unitriangular row substitutions in each letter
block, and is killed by
$\sum_{k=1}^n\partial_{u_{ik}}\partial_{v_{jk}}$ for every pair
of opposite rows at the same place. These conventions apply to
differential symbols as well as ordinary polynomial variables.

\begin{proposition}\label{arrep:reproduce}
Let $n$ be even, $2r\le n$, and let $P_1,\ldots,P_m$ be any finite
family of such polynomials, with at least one nonzero member. Denote by
$K(x,y)$ the result of applying
$\sum_j P_j(\partial_z)\overline{P_j}(\partial_w)$ at $z=w=0$
to the doubled Gaussian below, where the bar conjugates coefficients.
If $0<\delta\le1/(4(r+1))$, then there is a positive real number
$\rho_\delta$ such that every holomorphic $F$ on $\mathcal D$ with
$F(xb)=\chi_\lambda(b)F(x)$ and bounded absolute value on the
orthoframes satisfies the absolutely convergent identity
\[
         \int_{S_\delta} K(x,y)F(y)\,dy=\rho_\delta F(x)
                               \qquad(x\in\mathcal D).
\]
On orthoframes the pointwise automorphic norm in the bounded-family
construction equals $|F|$. The scalar depends on the specified measure,
archimedean data, family and slab. It applies to every deep level
using those same data.
\end{proposition}

\begin{proof}
First compute what survives the insertion into the Gaussian. Suppress
$\tau$, let $P_+,P_-$ be the diagonal projections to the two physical
blocks, and set $S=-J$, $\kappa=\tau(\kappa_0)$,
$h=\kappa/4$, $t=\pi i\kappa=-\pi\operatorname{Im}\kappa>0$. Put
\[
\begin{aligned}
 L&=P_+\bar y^tG^t+P_-x^{-1},&
 H&=P_+x^tG+P_-\bar y^{-1},\\
 W&=(L^{-1}SL)^tG,& Y&=(hW-(hW)^*)/(2i).
\end{aligned}
\]
The kernel of the bottom rows of $L$ is the positive space spanned
by the first $p$ columns of $x$. The kernel of its top rows is
the negative orthogonal complement of the first $p$ columns of $y$.
They intersect trivially. Interchanging $x,y$ and conjugating gives
invertibility of $H$. The involution $L^{-1}SL$ acts by $-1$ on the
first space and $+1$ on the second. For $v=v_++v_-$ in this
decomposition, the real part of $v^tW\bar v$ is
$-\langle v_+,v_+\rangle+\langle v_-,v_-\rangle$ because the
cross terms have zero real part. Thus $W+W^*$ is negative definite.
As $h/(2i)<0$, the Hermitian matrix $Y=(h/(2i))(W+W^*)$ is positive.

Write $M=L^{-t}GH^{-1}$ and $N=L^{-t}GY^{-1}GH^{-1}$.
For each column $k$ put $\alpha=P_+z_k+P_-w_k$,
$\beta=P_+w_k+P_-z_k$. The actual undifferentiated Gaussian has
prefactor $\prod_\tau(\det L\,\det H)^{-n/2}$ and exponent
$-4\pi\sum_{\tau,k}h_\tau^2\beta^t N\alpha$.
The off-diagonal blocks of $M^{-1}=HG^{-1}L^t$ vanish because
the corresponding products before taking blocks are identity
matrices. Its diagonal blocks are the leading $p$ block of
$x^tG\bar y$ and the trailing $q$ block of $\bar y^{-1}G^{-1}x^{-t}$.
In particular $SM=MS$. Moreover
\[
                     \frac{h}{2i}(SN+NS)=M .
\]
Indeed $W G^{-1}W=G$ and $W^*G^{-1}W^*=G$. For $T=W+W^*$
these give $WG^{-1}T=TG^{-1}W^*$ and therefore
$GT^{-1}W=W^*T^{-1}G$. It follows that
$WT^{-1}G+GT^{-1}W=G$. Multiply this last identity on the left
by $L^{-t}$ and on the right by $H^{-1}$,
using $L^{-t}W=SL^{-t}G$ and $WH^{-1}=GH^{-1}S$.
The latter equality follows from $SM=MS$.

The off-diagonal terms in $\sum_k\beta^t N\alpha$ are linear
combinations of $\sum_kz_{ik}z_{jk}$ and $\sum_kw_{ik}w_{jk}$
for opposite rows $i,j$. Pairing with a polynomial differential
symbol at zero makes multiplication by each such quadratic adjoint
to its mixed Laplacian. Every term containing one of these quadratics
vanishes, also after both sets of derivatives are taken. Only finitely
many degrees of the exponential enter the pairing. Thus one may
differentiate the following explicit expression instead:
\[
 \prod_\tau (\det L\,\det H)^{-n/2}
 \exp\!\left(\sum_{\tau,k}t_\tau
      [\,w_{+,k}^t M_{++}z_{+,k}-z_{-,k}^t M_{--}w_{-,k}\,]\right).
\]
It is holomorphic in $x$ and antiholomorphic in $y$.
The two diagonal blocks of $M(y,x)$ are the adjoints of those
of $M(x,y)$, and the prefactors interchange by complex conjugation.
Consequently $K(y,x)=\overline{K(x,y)}$.

Let $\mathcal U=\prod_\tau\{u:u^tG_\tau\bar u=G_\tau\}$. Under
simultaneous left translation by $u$, $L,H$ acquire right factors
$u^{-1},\bar u^{-1}$, so $M$ and the prefactor are invariant.
Under $x\mapsto xb$, write the two diagonal blocks of $b$ as $A,D$.
Then $L$ acquires left factor $1\oplus D^{-1}$ and $H$ left factor
$A^t\oplus1$. The prefactor gains $(\det D/\det A)^{n/2}$ at this
place. In the exponent the substitutions in $z$ are $A^{-t}$ on
the positive rows and $D$ on the negative rows. In the differential
symbol the resulting row substitutions are $A^{-1}$ and $D^t$.
Both are lower triangular in letter order. The prescribed row
degrees give exactly the remaining factor of $\chi_\lambda(b)$.
This proves equivariance in $x$, and symmetry gives the conjugate
character in $y$.

At $x=y=x_0$ one has $L=x_0^{-1}$, $H=\bar x_0^{-1}$ and $M=J$.
The exponential becomes $\exp(\sum_{\tau,i,k}t_\tau z_{\tau ik}
w_{\tau ik})$. With $P_{j,\nu}$ the monomial coefficients, ordinary
differentiation gives
\[
 C_0=K(x_0,x_0)=
   \left(\prod_\tau|\det x_{0,\tau}|^n\right)
   \sum_{j,\nu}|P_{j,\nu}|^2\,\nu!\,
                       \prod_\tau t_\tau^{|\nu_\tau|}>0 .
\]
Here $|\nu_\tau|$ is the sum of all exponents at that place and
$\nu!$ is the product of their factorials.

We next identify the only holomorphic function that can survive
the compact torus average. In coordinates $g=x_0^{-1}y$, let
$\Delta_{\tau,i}$ be the leading $i$-row principal minor of $g_\tau$.
Define
\[
 H_\lambda(y)=\prod_\tau
    \Delta_{\tau,r}^{\,\lambda_{\tau,r}}
    \prod_{i=1}^{r-1}\Delta_{\tau,i}^{
                              \lambda_{\tau,i}-\lambda_{\tau,i+1}} .
\]
For $r=0$ take the empty product. The weight decreases within each
block. When $0<p<r$ the only internal negative exponent occurs at
$i=p$, where the minor cannot vanish: the projection of a positive
$p$-plane to the standard positive coordinate space has trivial
kernel. Hence $H_\lambda$ is holomorphic throughout $\mathcal D$ and
has value $1$ at $x_0$. Under left diagonal unitary translation in
the $g$ coordinates it has character $\chi_\lambda$, and under right
triangular multiplication it has the same character.

These two transformations determine a holomorphic function up to its
value at $x_0$. Indeed on the big cell near $g=1$ write $g=\ell b$
with $\ell$ lower unitriangular. The two character laws imply
invariance of the restriction to $\ell$ under diagonal unitary
conjugation. On every strictly lower entry $\ell_{ij}$ the simultaneous
circle $\operatorname{diag}(e^{is},e^{2is},\ldots,e^{ris})$ acts with
positive weight $i-j$. The holomorphic power series of the restriction
therefore has only its constant term. Right equivariance determines
the function on a neighborhood, and the identity theorem extends the
conclusion. Indeed indefinite Gram--Schmidt identifies the domain with
the product of the orthoframes and the upper triangular matrices with
positive real diagonal. Both are connected: the unitary group modulo
$U(p)\times U(q)$ is the matrix ball described below, with connected
compact fiber. It follows that
\[
               K(y,x_0)=C_0H_\lambda(y).
\]
For this application the left torus law follows from simultaneous
invariance and right equivariance of $K$ in its second variable.

We check absolute integrability on the slab. The preceding
Gram--Schmidt coordinates write $y=x_0 g b$ with $g$ standard
$J$-unitary and $b$ upper triangular with positive real diagonal.
For our $\delta$, each leading Gram matrix has operator distance
at most $r\delta<1/4$ from the corresponding signed identity. Thus
even the closure of the permitted Gram matrices is compact with all
pivots nonzero and of the required signs. The signed triangular
factorization $b^tJ\bar b$ has a continuous inverse there by the
Gram--Schmidt recursion. The permitted $b$ therefore have compact
closure inside the triangular group. In these smooth coordinates
Lebesgue measure is $dg\,j(b)\,db$, where $dg$ is Haar measure and
$j$ is smooth positive: the original measure is left invariant
under $\mathcal U$ since $|\det u_\tau|=1$. Both $|\chi_\lambda(b)|$
and $j(b)$ are bounded on this compact closure. In particular
$F$ is bounded throughout $S_\delta$.

At a place with $pq>0$ write the standard unitary matrix
$g=\left(\begin{smallmatrix}A&B\\ C&D\end{smallmatrix}\right)$.
All singular values of $A$ and $D$ are at least one by the unitary
relations. The complementary-minor identity and $g^{-1}=Jg^*J$
give $|\det A|=|\det D|$ and bound each leading minor of size
$i>p$ by $|\det D|$ using the complementary trailing minor of
$g^{-1}$. For $i\le p$ the exterior powers of $A$ bound it by
$|\det A|$. As $|\det g|=1$, the product of the minor estimates,
keeping the exact denominator at $i=p$, gives the one-place bound
\[
       |H_{\lambda,\tau}(x_{0,\tau}g)|
          \le |\det A|^{\lambda_{\tau,1}-\lambda_{\tau,r}}
          \le |\det A|^{-n}\le1 .
\]
If either block is empty the group is compact and the function bounded.
For the mixed signature the group modulo the compact block diagonal
subgroup is the ball $Z=C A^{-1}$, $Z^*Z<I_p$, the graph coordinates
of its positive planes. The group is transitive by orthogonal
completion, and $|\det A|=\det(I_p-Z^*Z)^{-1/2}$.
The invariant density on the ball is proportional to
$\det(I_p-Z^*Z)^{-r}\,dZ$. To see this, for an acting matrix with
blocks $a,b,c,d$, set $U=a+bZ$ and $Z'=(c+dZ)U^{-1}$.
Then $I_p-Z'^*Z'=U^{-*}(I_p-Z^*Z)U^{-1}$ and the differential is
$H\mapsto(d-Z'b)HU^{-1}$. Block elimination gives
$|\det U|\,|\det(d-Z'b)|=|\det g_{\rm acting}|=1$, so the real
Jacobian is $|\det U|^{-2r}$. Hence integration of $|\det A|^{-n}$
on the group reduces, up to a finite constant, to integration of
$\det(I_p-Z^*Z)^{n/2-r}$ on a bounded ball. Its exponent is
nonnegative. This proves that $H_\lambda$ is bounded and absolutely
integrable on $S_\delta$, including at $n=2r$. In particular both
$K(x_0,y)F(y)$ and $|K(x_0,y)|^2$ are integrable there.

Finally let $t$ range over the product diagonal unitary torus with
probability Haar measure. Its action on frames is
$t\cdot y=x_0 t x_0^{-1}y$. The holomorphic projection
$F_\lambda(y)=\int\chi_\lambda(t)^{-1}F(t\cdot y)\,dt$ has the two
character laws above and has value $F(x_0)$ at $x_0$. Thus
$F_\lambda=F(x_0)H_\lambda$. The slab and its measure are invariant
under this torus. By absolute convergence and the conjugate torus
law of $K(x_0,y)$, replacing $F$ with $F_\lambda$ in the integral
at $x_0$ leaves its value unchanged. The resulting constant is
\[
                \rho_\delta=\frac{1}{C_0}
                   \int_{S_\delta}|K(x_0,y)|^2\,dy>0 .
\]
Positivity follows from continuity, $K(x_0,x_0)=C_0>0$, and the
open neighborhood of $x_0$ in the slab (with mass one in dimension
zero). Left translation carries the identity and absolute convergence
to every orthoframe by applying it to the translated bounded
holomorphic function. Right equivariance then carries it to every
frame of $\mathcal D$.
\end{proof}

The required invariant finite family can be chosen in these same
variables. At each place form the product of the initial $i$-row
minors to powers $c_i$ in the negative letters and to powers $c'_i$
in the positive letters, using respectively the first $q$ and the
last $p$ columns, nested within each set. Since $n\ge p+q$, these
sets are disjoint. The product over places is a nonzero row-highest
polynomial and all mixed Laplacians kill it. The entire space of
the stated harmonics lies in total degree
$d=\sum_{\tau,i}i(c_{\tau,i}+c'_{\tau,i})$, extending the arrays by
zero. Its dimension is at most $\binom{nr|\Psi|+d}{d}$, the number
of monomials of degree at most $d$. The compact column groups $U(n)^\Psi$ act on opposite letters
by a unitary matrix and its contragredient, preserving the row
conditions and the contractions. An invariant Hermitian product
exists by averaging. An orthonormal basis $P_1,\ldots,P_m$ then has
invariant sum $\sum_j P_j(z)\overline{P_j(w)}$ and positive $C_0$.
Thus the insertion used for spanning bounded forms satisfies the
proposition with its actual weight and with finite $m$.

\src{src/Development/AbelianVarietyTheory/Definitions/SpanningFamilies/IsSystem.lean}{65}{The doubled Gaussian before the two polynomial insertions};
\src{src/Development/AbelianVarietyTheory/Proofs/SiegelThetaTransformation/SetIntegralCompOfStable.lean}{268}{The positive reproducing scalar on the full frame slab}; and
\src{src/Development/AbelianVarietyTheory/Proofs/SiegelThetaTransformation/HermitianThetaSeriesSpanBoundedForms.lean}{168}{Spanning bounded families in the range $4r<n$}.
